\documentclass[10pt]{article}
\usepackage[utf8]{inputenc}
\usepackage[a4paper,height=20cm,width=13cm]{geometry}
\usepackage{amssymb}
\usepackage{dsfont}
\usepackage{calc}
\usepackage{graphicx}
% \usepackage{pst-node}
\usepackage{fourier}
\usepackage{euscript}
\usepackage{amsmath,amsthm}
\def\phi{\varphi}
\def\P{\EuScript P}
\def\M{\EuScript M}
\def\N{\EuScript N}
\def\B{\EuScript B}
\def\D{\EuScript D}
\def\C{\EuScript C}
\def\U{\EuScript U}
\def\V{\EuScript V}
\def\S{\EuScript S}
\def\X{\EuScript X}
\def\Y{\EuScript Y}
\def\Z{\EuScript Z}
\def\sm{\smallsetminus}
\def\niff{\nleftrightarrow}
\def\ZZ{\mathds Z}
\def\NN{\mathds N}
\def\PP{\mathds P}
\def\QQ{\mathds Q}
\def\RR{\mathds R}
\def\<{\langle}
\def\>{\rangle}
\def\E{\exists}
\def\A{\forall}
\def\0{\varnothing}
\def\imp{\rightarrow}
\def\iff{\leftrightarrow}
\def\IMP{\Rightarrow}
\def\IFF{\Leftrightarrow}
\def\range{\textrm{im}}
\def\Mod{\textrm{Mod}}
\def\Aut{\textrm{Aut}}
\def\Autf{\textrm{Aut$^{\rm f}$}}
\def\Th{\textrm{Th}}
\def\acl{\textrm{acl}}
\def\dcl{\textrm{dcl}}
\def\dom{{\rm dom}}
\def\eq{{\rm eq}}
\def\tp{\textrm{tp}}
\def\equivSh{\stackrel{\smash{\scalebox{.5}{\rm Sh}}}{\equiv}}
\def\equivL{\stackrel{\smash{\scalebox{.5}{\rm L}}}{\equiv}}

\def\cnonfork{\mathbin{\raise1.8ex\rlap{\kern0.6ex\rule{0.6ex}{0.1ex}}\rlap{\kern1.1ex\rule{0.1ex}{1.9ex}}\raise-0.3ex\hbox{$\smile$} } }

\newcommand{\labella}[1]{{\sf\footnotesize #1}\hfill}
\renewenvironment{itemize}
  {\begin{list}{$\triangleright$}{%
   \setlength{\parskip}{0mm}
   \setlength{\topsep}{0mm}
   \setlength{\rightmargin}{0mm}
   \setlength{\listparindent}{0mm}
   \setlength{\itemindent}{0mm}
   \setlength{\labelwidth}{3ex}
   \setlength{\itemsep}{0mm}
   \setlength{\parsep}{0mm}
   \setlength{\partopsep}{0mm}
   \setlength{\labelsep}{1ex}
   \setlength{\leftmargin}{\labelwidth+\labelsep}
   \let\makelabel\labella}}{%
   \vspace*{-.3\baselineskip}
  \end{list}}
%\def\ssf#1{\textsf{\small #1}}
\newcounter{ex}
\newenvironment{exercise}{\bigskip\addtocounter{ex}{1}\textbf{Esercizio \theex.\ }}{}
\pagestyle{empty}
\parindent0ex
\parskip1ex
\raggedbottom
\def\nsR{{}^*\!\RR}
\def\ssf#1{\textsf{#1}}
\renewcommand{\baselinestretch}{1.3}


\usepackage{fancyhdr}
\pagestyle{fancy}
\lfoot{Teoria dei Modelli a.a.~2025/26}
\rhead{}
\rfoot{{\small\thepage}}
\cfoot{}

%
% INDICARE NOME E COGNOME DI TUTTI GLI AUTORI
%
\lhead{Nome/i Cognome/i}
%

\begin{document}

\begin{exercise}
  (Alessandro Martina)
  Prove that $G$-syndetic sets are weakly $G$-thick.
\end{exercise}

\begin{exercise} 
  (Francesco Sulpizi)
  Assume that $\X$ is $\Delta(\Z)$-type-definable and $G$ acts transitively on $\X$.
  Let $\Y\subseteq\X$ be a $\Delta(\Z)$-type-definable set with a small $G$-orbit.
  Prove that $\Y$ is definable by a $G$-syndetic type.
\end{exercise}

\begin{exercise}
  (Leonardo Centazzo)
  Prove that the following are equivalent
  \begin{itemize}
    \item[1.] $\D$ is weakly $G$-thick
    \item[2.] $\D=\C\cap\B$ for some $G$-syndetic set $\B$ and some $G$-thick set $\C$
    \item [3.] there is a non $G$-syndetic set $\C$ such that $\D\cup\C$ is syndetic.
  \end{itemize}
\end{exercise}

\begin{exercise}
  (Roberto Carnevale)
\end{exercise}

\begin{exercise}
  (Matteo Bisi)
  Let $\D$ be a $\Delta(\Z)$-type-definable set.
  Are the following are equivalent?
  \begin{itemize}
    \item[1.] $\D$ is $G$-thick
    \item[2.] $\D$ is definable by a $G$-thick type.
  \end{itemize}
\end{exercise}

\begin{exercise}
  (Davide Peccioli)
\end{exercise}

\begin{exercise}
  (Pietro Giura)
  Let $\X=\Z=\U$.
  Prove that if $p(x)\in S(\U)$ is finitely satisfiable in every $M\supseteq A$ then it is thick under the action of $G=\Autf(\U)$.
  Is the same true for incomplete types?
\end{exercise}

\begin{exercise} 
  (Costanza Furone)  
  Prove that $G$-syndetic sets are weakly $G$-thick.
\end{exercise}

\begin{exercise}
  (Angela Dosio)
\end{exercise}
  
\end{document}


